% ==========================================================================
%  Chapter 49 — Phase 4: Finite-Strain Kinematics and Hysteretic Cycling
% ==========================================================================

\chapter{Phase 4: Finite-Strain Kinematics and Hysteretic Cycling}
\label{ch:phase4-finite-strain-hysteretic}

This chapter documents Phase~4 of the master plan
(Chapter~\ref{ch:convergence-master-plan}).
The original scope required implementing the Updated Lagrangian (UL)
framework: deformation gradient computation, spatial tangent evaluation,
geometric stiffness assembly, and reference configuration update.
A full codebase audit revealed that the finite-strain kinematic
infrastructure was \emph{already completely implemented and tested},
covering Total Lagrangian (TL), Updated Lagrangian (UL), and all
required strain/stress measures.

Phase~4 was therefore extended to include a comprehensive verification
of the cyclic constitutive models already present in the library—%
Menegotto--Pinto steel, Kent--Park concrete, and composable J$_2$
plasticity—under prescribed cyclic strain protocols, with CSV output
and automated stress--strain hysteresis plotting.


%% ------------------------------------------------------------------
\section{Phase 4 Scope Revision}
\label{sec:phase4-scope-revision}

Chapter~\ref{ch:convergence-master-plan} defined Phase~4 as
``Finite-Strain Kinematics: Updated Lagrangian (UL) framework'' with
four deliverables:

\begin{enumerate}
  \item Deformation gradient
        $\mathbf{F} = \mathbf{I} + \nabla\mathbf{u}$ from shape
        function gradients.\label{it:p4-F}
  \item Spatial tangent modulus $\mathbb{c}$ via the constitutive
        interface.\label{it:p4-c}
  \item Geometric stiffness
        $K_\sigma$ alongside material tangent $K_m$.\label{it:p4-Ks}
  \item Reference configuration update after each converged
        step.\label{it:p4-ref}
\end{enumerate}

After auditing \texttt{src/continuum/}, all four deliverables were
found to be \textbf{operational}, with extensive test coverage:

\begin{center}
\begin{tabular}{lcc}
  \toprule
  Test suite & Assertions & Status \\
  \midrule
  \texttt{test\_kinematic\_policy} (TL + UL traits, B matrices)
      & 68 & \checkmark \\
  \texttt{test\_hyperelastic} (SVK, Neo-Hookean, energy/tangent)
      & 69 & \checkmark \\
  \texttt{test\_tl\_integration} (element-level TL pipeline)
      & 48 & \checkmark \\
  \texttt{test\_updated\_lagrangian} (UL $\equiv$ TL equivalence)
      & 41 & \checkmark \\
  \texttt{test\_finite\_strain\_damage} (damage + hyperelastic)
      & 25 & \checkmark \\
  \midrule
  Total pre-existing & \textbf{251} & all pass \\
  \bottomrule
\end{tabular}
\end{center}

Given this, Phase~4 was extended with a second objective: comprehensive
cyclic constitutive model verification with hysteretic stress--strain
curve output.


%% ------------------------------------------------------------------
\section{Finite-Strain Infrastructure Summary}
\label{sec:phase4-finite-strain}

\subsection{Kinematic Policies}

The library provides four compile-time kinematic policies in
\texttt{KinematicPolicy.hh}, each satisfying the
\texttt{KinematicPolicyConcept}:

\begin{center}
\begin{tabular}{lccc}
  \toprule
  Policy & Geometrically linear & $K_\sigma$ & Current-vol factor \\
  \midrule
  \texttt{SmallStrain}       & \checkmark & --- & --- \\
  \texttt{TotalLagrangian}   & ---        & \checkmark & --- \\
  \texttt{UpdatedLagrangian} & ---        & \checkmark & \checkmark \\
  \texttt{Corotational}      & ---        & \checkmark & placeholder \\
  \bottomrule
\end{tabular}
\end{center}

\subsection{Deformation Gradient}

\texttt{TotalLagrangian::compute\_F\_from\_gradients<dim>()} computes
\[
  F_{ij} = \delta_{ij} + \sum_I u_{iI}\,\frac{\partial N_I}{\partial X_j}
\]
from the shape function gradient matrix $\nabla\mathbf{N}$ and the flat
displacement vector $\mathbf{u}_e$.  The Updated Lagrangian path reuses
the same $F$ computation and additionally evaluates spatial gradients
$\nabla_x N = \nabla_X N \cdot F^{-1}$ for the spatial B operator and
geometric stiffness.

\subsection{Strain and Stress Measures}

\texttt{StrainMeasures.hh} and \texttt{StressMeasures.hh} provide the
full tensor algebra toolkit:

\begin{itemize}
  \item Green--Lagrange
        $\mathbf{E} = \tfrac{1}{2}(\mathbf{F}^T\mathbf{F} - \mathbf{I})$
  \item Almansi
        $\mathbf{e} = \tfrac{1}{2}(\mathbf{I} - \mathbf{b}^{-1})$
  \item Hencky (logarithmic) strain
  \item Seth--Hill family (parametric)
  \item Piola transformations:
        $\mathbf{S} \to \boldsymbol{\sigma}$ (push-forward),
        $\boldsymbol{\sigma} \to \mathbf{S}$ (pull-back),
        Kirchhoff $\boldsymbol{\tau} = J\boldsymbol{\sigma}$,
        first Piola--Kirchhoff $\mathbf{P} = \mathbf{F}\mathbf{S}$
\end{itemize}

\subsection{Hyperelastic Models}

Two models satisfy \texttt{HyperelasticModelConcept}:

\begin{itemize}
  \item \texttt{SaintVenantKirchhoff<dim>}: $W = \tfrac{\lambda}{2}(\mathrm{tr}\,E)^2 + \mu\,E : E$
  \item \texttt{CompressibleNeoHookean<dim>}: $W = \tfrac{\mu}{2}(\mathrm{tr}\,C - d) - \mu\ln J + \tfrac{\lambda}{2}(\ln J)^2$
\end{itemize}

Both provide $W(E)$, $S(E)$, and $\mathbb{C}(E)$.  The
\texttt{HyperelasticRelation} adapter bridges the tensor API to the
general \texttt{ConstitutiveRelation} interface (Voigt convention).

\subsection{Geometric Stiffness}

Both TL and UL paths assemble the geometric stiffness contribution:

\begin{itemize}
  \item TL: $(K_\sigma)_{IiJj} = \delta_{ij}\int_{\Omega_0} S_{kl}\,
        \frac{\partial N_I}{\partial X_k}\,
        \frac{\partial N_J}{\partial X_l}\,dV$
  \item UL: $(k_\sigma)_{IiJj} = \delta_{ij}\int_{\Omega_n} \sigma_{kl}\,
        \frac{\partial N_I}{\partial x_k}\,
        \frac{\partial N_J}{\partial x_l}\,dv$
\end{itemize}

\texttt{ContinuumElement} dispatches via \texttt{if constexpr} on
\texttt{KinematicPolicy::needs\_geometric\_stiffness} and
\texttt{KinematicPolicy::needs\_current\_volume\_factor} to assemble
$K_t = K_m + K_\sigma$ correctly.


%% ------------------------------------------------------------------
\section{Cyclic Constitutive Model Verification}
\label{sec:phase4-hysteretic-cycling}

Three cyclic/hysteretic material models were exercised under prescribed
strain protocols:

\subsection{Menegotto--Pinto Steel}

\texttt{MenegottoPintoSteel} implements the Menegotto \& Pinto (1973)
uniaxial model with the Filippou et~al.\ (1983) modifications.
Key features: Bauschinger effect via evolving curvature parameter $R$,
isotropic hardening amplitude, and yield plateau.

\noindent
Constructor: \texttt{MenegottoPintoSteel(E0, fy, b, R0, cR1, cR2, a1, a2)}.

\subsubsection*{Test 1: Increasing-Amplitude Protocol}

A seismic-standard protocol $(\pm 1\varepsilon_y, \pm 2\varepsilon_y,
\pm 4\varepsilon_y, \pm 8\varepsilon_y, \pm 12\varepsilon_y)$ with
30~subdivisions per segment (301 points total).  Assertions:

\begin{itemize}
  \item Peak stresses $\geq 0.95 f_y$ in both tension and compression
  \item Bauschinger effect detected: $|\sigma| < f_y$ at
        $\varepsilon = -\varepsilon_y$ during first compression reversal
  \item Non-zero hysteretic loop area (positive energy dissipation)
\end{itemize}

\subsubsection*{Test 2: Constant-Amplitude (3 Cycles)}

Three full cycles at $\pm 4\varepsilon_y$ with 40~subdivisions per
segment.  Validates that each cycle dissipates positive absolute work,
and that the R-evolution stabilizes the loop shape.

\subsection{Kent--Park Concrete}

\texttt{KentParkConcrete} implements the Scott--Park--Priestley (1982)
model for confined and unconfined concrete.  Features: parabolic
ascending branch, linear softening, origin-pointing unloading, tension
cutoff.

\subsubsection*{Test 3: Unconfined Cyclic}

Protocol: $0 \to -0.001 \to 0 \to -0.003 \to +0.0002 \to -0.005 \to 0$.
Assertions:

\begin{itemize}
  \item Peak compression $\geq 80\%$ of $f'_c$
  \item Origin-pointing unloading: stress $\approx 0$ at $\varepsilon=0$
  \item Tension stress $\leq f_t + 1$\,MPa
  \item Compression softening at $\varepsilon = -0.005$
\end{itemize}

\subsubsection*{Test 4: Confined Cyclic}

Same concrete with confinement parameters
($\rho_s = 0.015$, $f_{yh} = 400$\,MPa, $h' = 300$\,mm, $s_h = 100$\,mm).
Protocol extends to $\varepsilon = -0.015$.  Validates:

\begin{itemize}
  \item Confined peak stress $> f'_c$ (confinement factor $K > 1$)
  \item Residual stress $> 10\%\,f'_c$ at $\varepsilon = -0.015$
\end{itemize}

\subsection{J$_2$ Plasticity (Uniaxial)}

\texttt{J2PlasticityRelation<UniaxialMaterial>} with linear isotropic
hardening ($E = 200\,000$, $\sigma_y = 250$, $H = 5000$).

\subsubsection*{Test 5: Symmetric Cyclic (3 Cycles at $\pm 3\varepsilon_y$)}

Validates:
\begin{itemize}
  \item Peak stress exceeds initial yield (isotropic hardening)
  \item Symmetric response: $|\sigma_{\max} + \sigma_{\min}| < 5\%\,\sigma_{\max}$
  \item Non-zero hysteretic loop area
\end{itemize}

\subsection{Commit/Revert Protocol}

\subsubsection*{Test 6: Idempotency Under Cycling}

After a full plastic cycle, a trial strain is computed \emph{without
committing}.  Verifies that the response at a reference strain is
unchanged (the \texttt{ExternallyStateDriven} interface is read-only
on committed state).  After committing the trial, the response changes
(path dependence confirmed).

\subsection{Combined Output}

\subsubsection*{Test 7: Steel vs.\ Concrete on Same Protocol}

Both materials driven through the same $\pm 0.015$ strain protocol.
A combined CSV is written with columns \texttt{strain},
\texttt{steel\_stress}, \texttt{concrete\_stress} for direct comparison.

\subsection{Finite-Strain Gate}

\subsubsection*{Test 8: Kinematic Policy Verification}

Compile-time trait checks for all three policies, plus runtime
deformation gradient computation at zero and small displacement (verifies
$F = I$ and $F \approx I + \nabla u$).


%% ------------------------------------------------------------------
\section{CSV Output and Hysteresis Plotting}
\label{sec:phase4-csv-plotting}

All cyclic tests write CSV files to
\texttt{data/output/hysteretic\_cycles/}:

\begin{center}
\begin{tabular}{ll}
  \toprule
  File & Content \\
  \midrule
  \texttt{steel\_menegotto\_pinto\_cyclic.csv}         & Increasing-amplitude steel \\
  \texttt{steel\_menegotto\_pinto\_3cycles.csv}         & Constant-amplitude 3 cycles \\
  \texttt{concrete\_kent\_park\_cyclic.csv}             & Unconfined concrete \\
  \texttt{concrete\_kent\_park\_confined\_cyclic.csv}   & Confined concrete \\
  \texttt{j2\_uniaxial\_cyclic.csv}                    & J$_2$ plasticity \\
  \texttt{combined\_steel\_concrete\_cyclic.csv}        & Steel + concrete comparison \\
  \bottomrule
\end{tabular}
\end{center}

A Python~3 script \texttt{scripts/plot\_hysteresis\_curves.py} reads
these CSVs and generates a multi-panel figure with matplotlib
(\texttt{hysteresis\_curves.pdf} and \texttt{.png}).  Each panel shows
the stress--strain hysteresis loop for one material/protocol, with the
combined panel overlaying steel and concrete on the same axes.


%% ------------------------------------------------------------------
\section{Results}
\label{sec:phase4-results}

\begin{center}
\begin{tabular}{lcc}
  \toprule
  Metric & Target & Actual \\
  \midrule
  Phase~4 hysteretic cycling assertions & 41 & \textbf{41/41} \\
  Pre-existing finite-strain tests & 251 & \textbf{251/251} \\
  Full ctest suite & $\geq 43/45$ & \textbf{43/45} \\
  Pre-existing failures & 2 & 2 \\
  New regressions & 0 & \textbf{0} \\
  Build errors (\texttt{-Werror}) & 0 & \textbf{0} \\
  \bottomrule
\end{tabular}
\end{center}


%% ------------------------------------------------------------------
\section{Phase 4 Completion Status}
\label{sec:phase4-completion}

\begin{center}
\begin{tabular}{lcc}
  \toprule
  Deliverable & Status & Notes \\
  \midrule
  Deformation gradient $\mathbf{F}$ & \checkmark & Pre-existing (TL + UL) \\
  Spatial tangent $\mathbb{c}$ & \checkmark & Pre-existing (push-forward) \\
  Geometric stiffness $K_\sigma$ & \checkmark & Pre-existing (TL + UL) \\
  Reference config update & \checkmark & Pre-existing (UL spatial grads) \\
  Menegotto--Pinto cycling & \checkmark & Phase~4 test \\
  Kent--Park cycling & \checkmark & Phase~4 test \\
  J$_2$ uniaxial cycling & \checkmark & Phase~4 test \\
  Commit/revert idempotency & \checkmark & Phase~4 test \\
  CSV curve output & \checkmark & Phase~4 test \\
  Python hysteresis plotter & \checkmark & Phase~4 script \\
  \midrule
  Corotational formulation & --- & Deferred (placeholder only) \\
  Kinematic hardening (Chaboche) & --- & Deferred (future extension) \\
  \bottomrule
\end{tabular}
\end{center}

\noindent
Phase~4 is complete.  The finite-strain kinematic framework was found
fully operational, and the cyclic constitutive models have been verified
under aggressive cyclic protocols with 41 assertions, CSV output, and
automated stress--strain hysteresis plotting.
